%\newpage
\section{The Fourier transform of the Faraday tensor - direct formula}\label{sect:FTF-direct}
The goal is to calculate the Fourier transform of the Faraday tensor
\begin{align}
    F(\omega, {\bf x}) = \frac1{2\pi}\int\limits_{-\infty}^{\infty} dt e^{i\omega t}F(t,{\bf x})
\end{align}
by the direct transformation of the primary form
\begin{align}
    F^{\alpha\beta}(x)
    =\left.\frac{e}{4\pi\epsilon_0c}
    \frac{(u\cdot R_0)(R_0^{\alpha}w^{\beta}-R_0^{\beta}w^{\alpha})
        -((w\cdot R_0)-c^2)(R_0^{\alpha}u^{\beta}-R_0^{\beta}u^{\alpha})}
    {(u\cdot R_0)^3}\right|_{\tau=\tau_r}.\label{eq:F_alpha_beta_primary_form-used1}
\end{align}
The fourier transform becomes
\begin{align}
    F^{\alpha\beta}(\omega, {\bf x})=\frac1{2\pi}\int\limits_{-\infty}^{\infty}dte^{i\omega t}F(t,{\bf x})
\end{align}
we make the change of variable from $t$ to $\tau_r$. We have
The relation between the laboratory time $t$ and the retarded proper time $\tau_r$ is given in (\ref{ct-taur-rel})
\begin{align}
    ct= r^0(\tau_r)+|{\bf x}-{\bf r}(\tau_r)| \equiv r_0^0(\tau_r)+|{\bf x}_0-{\bf r}_0(\tau_r)|.
\end{align}
We change tha variable of integration for $t$ to $\tau_r$ in the previous integral, using
\begin{align}
    \frac{d{ct}}{d\tau_r} = \frac{dr_0^0(\tau_r)}{d\tau_r} + \frac{d|{\bf x}_0 -{\bf r}_0(\tau_r)|}{d\tau_r} = u^0(\tau_r) - {\bf u}(\tau_t)\cdot\frac{{\bf x}_0 -{\bf r}_0(\tau_r)}{{|\bf x}_0 -{\bf r}_0(\tau_r)|}
\end{align}
or, using the unity four vectors $R_0$ and   $n_{R_0}$ defined in (\ref{defnr})
\begin{align}
    n_{R_0}=(1,{\bf n}_{R_0}),\quad {\bf n}_{R_0}=\frac{{\bf x}_0-{\bf r}_0(\tau_r)}{|{\bf x}_0-{\bf r}_0(\tau_r)|},\quad  R_0 =  |{\bf R}_0|n_{R_0}
\end{align}
we get
\begin{importantbox}\begin{align}
        \frac{d(ct)}{d\tau_r} = {n_{R_0}(\tau_r)\cdot u(\tau_r)}
    \end{align}\end{importantbox}
Note that this is the derivative of $ct$, not of $t$: changing the integration
variable in $\int dt\, e^{i\omega t}(\ldots)$ requires the dimensionless
Jacobian
\begin{importantbox}
    \begin{align}
        \frac{dt}{d\tau_r}=\frac1c\, n_{R_0}(\tau_r)\cdot u(\tau_r)\label{dtdtaur}
    \end{align}
\end{importantbox}
which must be inserted explicitly, and not dropped, when the integration
variable is changed from $t$ to $\tau_r$.
Using the notation
\begin{importantbox}
    \begin{align}
        k = \frac{\omega}c\label{defk-primar}
    \end{align}
\end{importantbox}
and denoting the variable of integration by $\tau$, we have
\begin{align}
    F^{\alpha\beta}(\omega,{\bf x}) & =\frac1{2\pi}\int\limits_{-\infty}^{\infty}\frac{dt}{d\tau}\,d\tau\, e^{ik\left(r_0^0(\tau)+|{\bf x}_0-{\bf r}_0(\tau)|\right)}F^{\alpha\beta}(\tau)\nonumber                                                                                                            \\
                                    & =\frac1{2\pi}\frac{e}{4\pi\epsilon_0c^2}\int\limits_{-\infty}^{\infty}d\tau\,(u\cdot n_{R_0})\, e^{ik\left(r_0^0(\tau)+|{\bf x}_0-{\bf r}_0(\tau)|\right)}\frac{(u\cdot R_0)(R_0^{\alpha}w^{\beta}-R_0^{\beta}w^{\alpha})
        -((w\cdot R_0)-c^2)(R_0^{\alpha}u^{\beta}-R_0^{\beta}u^{\alpha})}
    {(u\cdot R_0)^3}
\end{align}
using $u\cdot n_{R_0}=(u\cdot R_0)/|{\bf R}_0|$ (\ref{defnr}), one power of
$(u\cdot R_0)$ cancels and this simplifies to
\begin{importantbox}
    \begin{align}
        F^{\alpha\beta}(ck,{\bf x})=\frac1{2\pi}\frac{e}{4\pi\epsilon_0c^2|{\bf R}_0|}\int\limits_{-\infty}^{\infty}d\tau e^{ik\left(r_0^0(\tau)+|{\bf x}_0-{\bf r}_0(\tau)|\right)}\frac{(u\cdot R_0)(R_0^{\alpha}w^{\beta}-R_0^{\beta}w^{\alpha})
        -((w\cdot R_0)-c^2)(R_0^{\alpha}u^{\beta}-R_0^{\beta}u^{\alpha})}
    {(u\cdot R_0)^2}\label{eq:F-direct-corrected}
    \end{align}
\end{importantbox}
In the previous result we can isolate the long from the short range components
\begin{align}
    F_l^{\alpha\beta}(ck,{\bf x}) & =\frac1{2\pi}\frac{e}{4\pi\epsilon_0c^2|{\bf R}_0|}\int\limits_{-\infty}^{\infty}d\tau e^{ik\left(r_0^0+|{\bf R}_0|\right)}\frac{(u\cdot R_0)(R_0^{\alpha}w^{\beta}-R_0^{\beta}w^{\alpha})-(w\cdot R_0)(R_0^{\alpha}u^{\beta}-R_0^{\beta}u^{\alpha})}
    {(u\cdot R_0)^2}\nonumber                                                                                                                                                                                                                                                \\
                                  & =\frac1{2\pi}\frac{e}{4\pi\epsilon_0c^2|{\bf R}_0|}\int\limits_{-\infty}^{\infty}d\tau e^{ik\left(r_0^0+|{\bf R}_0\right)}\frac{(u\cdot n_{R_0})(n_{R_0}^{\alpha}w^{\beta}-n_{R_0}^{\beta}w^{\alpha})-(w\cdot n_{R_0})(n_{R_0}^{\alpha}u^{\beta}-n_{R_0}^{\beta}u^{\alpha})}
    {(u\cdot n_{R_0})^2}
\end{align}
and similarly
\begin{align}
    F_s^{\alpha\beta}(ck,{\bf x})=\frac1{2\pi}\frac{e c^2}{4\pi\epsilon_0c^2|{\bf R}_0|}\int\limits_{-\infty}^{\infty}d\tau e^{ik\left(r_0^0+|{\bf R}_0|\right)}\frac{(R_0^{\alpha}u^{\beta}-R_0^{\beta}u^{\alpha})}
    {(u\cdot R_0)^2}=\frac1{2\pi}\frac{e}{4\pi\epsilon_0|{\bf R}_0|^2}\int\limits_{-\infty}^{\infty}d\tau e^{ik\left(r_0^0+|{\bf R}_0|\right)}\frac{(n_{R_0}^{\alpha}u^{\beta}-n_{R_0}^{\beta}u^{\alpha})}
    {(u\cdot n_{R_0})^2}
\end{align}

The final results are
\begin{highlightbox}
    \begin{align}
        F^{\alpha\beta}(c k,{\bf x})  & = F^{\alpha\beta}_{l}(c k,{\bf x}) + F^{\alpha\beta}_{s}(c k,{\bf x}),
        \\
        F_l^{\alpha\beta}(ck,{\bf x}) & =\frac1{2\pi}\frac{e}{4\pi\epsilon_0c^2|{\bf R}_0|}\int\limits_{-\infty}^{\infty}d\tau e^{ik\left(r_0^0+|{\bf R}_0\right)}\frac{(u\cdot n_{R_0})(n_{R_0}^{\alpha}w^{\beta}-n_{R_0}^{\beta}w^{\alpha})-(w\cdot n_{R_0})(n_{R_0}^{\alpha}u^{\beta}-n_{R_0}^{\beta}u^{\alpha})}
        {(u\cdot n_{R_0})^2}                                                                                                                                                                                                                                             \\
        F_s^{\alpha\beta}(ck,{\bf x}) & =\frac1{2\pi}\frac{e}{4\pi\epsilon_0|{\bf R}_0|^2}\int\limits_{-\infty}^{\infty}d\tau e^{ik\left(r_0^0+|{\bf R}_0|\right)}\frac{(n_{R_0}^{\alpha}u^{\beta}-n_{R_0}^{\beta}u^{\alpha})}
        {(u\cdot n_{R_0})^2}
    \end{align}
\end{highlightbox}
Compare with the simplified form of Sec.~\ref{sect:FTF-simplified}: both
derivations now carry the same overall factor $e/(4\pi\epsilon_0c^2)$, as
required for the two routes to compute the same $F^{\alpha\beta}(ck,{\bf x})$.
The individual $F_l$, $F_s$ pieces need not agree term by term, since the
long/short split is defined differently in each derivation; only the total
$F^{\alpha\beta}=F_l^{\alpha\beta}+F_s^{\alpha\beta}$ is guaranteed to match.
The notations used here are those already defined
\begin{highlightbox}
    \begin{align}
        k = \frac{\omega}{c},\quad
        R_0 = (|{\bf R}_0|,{\bf R}_0)=|{\bf R}_0| (1,{\bf n}_{R_0}) = |{\bf R}_0|n_{R_0}, \qquad {\bf R}_0 = {\bf x}_0-{\bf r}_0(\tau),\quad {\bf n}_{R_0} = \frac{{\bf R}_0}{|{\bf R}_0|}.
    \end{align}
\end{highlightbox}
